% !TEX root = Master_PS.tex
%=======================================
% 02 SAFETY NETS -- family + government support, my brother's unemployment
% Theme: how my personal background motivated a graduate degree
% (also covers opportunities/challenges: Michigan, Wisconsin, Berkeley, Northwestern)
%=======================================
As I look back on my high school and university experiences, I think a lot about how important my family and the informal safety net they provided was to my academic and professional journey. I could take risks in my career, take on more study or choose more demanding professional roles in large part because I had informal support. But I also think a lot about how getting through the door to these opportunities required government support; I went to a public university, used government-issued university loans, and relied income assistance when studying full-time. In this way, both government assistance and family support played complementary roles in helping my career. But I've also seen examples where both fall short. I watched as my brother was unemployed for over three years and needed more than just the imperfect income assistance provided by unemployment insurance. Like most young Australians, he is college educated, but job-search in fields like IT and computer science is incredibily challenging because these jobs are rapidly changing to due AI. Government programs aren't well-positioned to alleviate labour-market scarring in the face of such technological disruption. At the same time, moral support from family can often be viewed as patronising or insensitive, and financial support to wedge the gap that unemployment insurance creates can only last so long.

% Note to self: the informal safety net my parents provided had important
% spillover effects on their own lives -- effectively their own lives had to be
% put on hold. How important the informal safety net is in informing individual
% decisions.
